\documentclass[12pt,a4paper]{article}
\usepackage[utf8]{inputenc}
\usepackage[T1]{fontenc}
\usepackage[brazilian]{babel}
\usepackage{geometry}
\usepackage{booktabs}
\usepackage{array}
\usepackage{hyperref}
\usepackage{enumitem}

\geometry{margin=2.5cm}

\title{
    \textbf{Segurança da Informação} \\
    \large Avaliação N3, Checkpoint de Fundação Técnica (09/06/2026) \\
    \normalsize Curso de Engenharia de Software
}
\author{Grupo Lendify}
\date{9 de junho de 2026}

\begin{document}

\maketitle

\section*{Identificação do Projeto}

\begin{tabular}{@{}ll@{}}
\toprule
\textbf{Código do Projeto} & P03-B, Reserva de Equipamentos \\
\textbf{Nome do Sistema}   & Lendify \\
\textbf{Professor}         & Edson Vaz Lopes \\
\textbf{Data}              & 09 de junho de 2026 \\
\bottomrule
\end{tabular}

\section*{Integrantes do Grupo}

\begin{enumerate}
    \item Miguel Angelo Dufloth Filho
    \item Miguel Angel Balladares Huertas
    \item Luis Fernando Pereira
    \item Leonardo Lotério de Lima
    \item Lucas Honorato dos Santos
\end{enumerate}

\section*{Repositório GitHub}

Link: \url{https://github.com/migueldufloth/lendify} \\
Acesso do professor: confirmado.

\section*{Descrição Curta do Sistema}

O \textbf{Lendify} é um sistema web para gerenciamento de empréstimo e reserva de
equipamentos como projetores, notebooks e kits. Solicitantes realizam pedidos de
reserva, operadores aprovam e acompanham devoluções, e administradores gerenciam
o acervo, os usuários e os logs de auditoria.

O nome é uma combinação de \textit{lend} (emprestar, em inglês) com o sufixo
\textit{-fy}, amplamente utilizado em nomenclatura de produtos tecnológicos.

\section*{Stack Pretendida}

\begin{itemize}
    \item \textbf{Front-end:} React
    \item \textbf{Back-end:} Django (Python)
    \item \textbf{Banco de dados:} MySQL
\end{itemize}

\section*{Perfis de Usuário Previstos}

\begin{tabular}{@{}p{3.5cm}p{10cm}@{}}
\toprule
\textbf{Perfil} & \textbf{Responsabilidade} \\
\midrule
Solicitante   & Realiza pedidos de reserva e visualiza apenas suas próprias solicitações. \\
Operador      & Analisa, aprova ou recusa reservas e registra devoluções. \\
Administrador & Gerencia equipamentos, usuários, permissões e logs de auditoria. \\
\bottomrule
\end{tabular}

\section*{Funcionalidades Mínimas Previstas}

\subsection*{Solicitante}
\begin{itemize}
    \item Realizar reserva ou empréstimo de equipamento.
    \item Visualizar suas próprias reservas.
    \item Cancelar reserva com status pendente.
\end{itemize}

\subsection*{Operador}
\begin{itemize}
    \item Visualizar reservas sob sua responsabilidade.
    \item Aprovar ou recusar solicitações.
    \item Registrar devolução de equipamento.
\end{itemize}

\subsection*{Administrador}
\begin{itemize}
    \item Cadastrar, editar e remover equipamentos.
    \item Gerenciar usuários e perfis de acesso.
    \item Visualizar logs de auditoria.
    \item Gerar relatório de uso dos equipamentos.
\end{itemize}

\section*{Arquitetura Inicial}

O Lendify adotará uma arquitetura \textit{Client-Server} baseada em APIs REST, focada em separação de responsabilidades e segurança de acesso:

\begin{itemize}
    \item \textbf{Front-end (Client):} Desenvolvido em React. Consumirá a API de forma assíncrona, enviando os tokens de autenticação (JWT) no cabeçalho das requisições para atestar a identidade e perfil do usuário.
    \item \textbf{Back-end (API):} Desenvolvido em Django utilizando o Django REST Framework. Atuará como a camada central de validação, aplicando as regras de negócio de reservas, protegendo as rotas por perfil e gravando os logs.
    \item \textbf{Banco de Dados:} MySQL, acessado estritamente através do ORM do Django, evitando exposição direta a SQL Injections.
\end{itemize}

\section*{Matriz Inicial de Permissões}

\begin{tabular}{@{}lccc@{}}
\toprule
\textbf{Funcionalidade} & \textbf{Solicitante} & \textbf{Operador} & \textbf{Administrador} \\
\midrule
Login no sistema & Sim & Sim & Sim \\
Visualizar equipamentos & Sim & Sim & Sim \\
Gerenciar equipamentos (CRUD) & Não & Não & Sim \\
Solicitar reserva & Sim & Não & Não \\
Visualizar próprias reservas & Sim & Sim & Sim \\
Visualizar todas as reservas & Não & Sim & Sim \\
Aprovar/Recusar reservas & Não & Sim & Sim \\
Registrar retirada/devolução & Não & Sim & Sim \\
Visualizar logs de auditoria & Não & Não & Sim \\
Gerenciar usuários & Não & Não & Sim \\
\bottomrule
\end{tabular}

\section*{Modelo Inicial de Dados (Entidades Principais)}

\begin{itemize}
    \item \textbf{User (Auth):} Armazena credenciais, hash de senha e perfil de acesso.
    \item \textbf{Equipamento:} Representa o patrimônio. Contém nome, descrição, número de série e status (ativo/inativo).
    \item \textbf{Reserva:} A entidade central do fluxo de empréstimos. Possui chaves estrangeiras para \textit{User} e \textit{Equipamento}. Armazena data/hora de início, data/hora de fim, justificativas e o \textbf{status} atual (Pendente, Aprovada, Recusada, Em Andamento, Devolvida, Cancelada).
    \item \textbf{LogAuditoria:} Registra as ações de mudança de estado sensíveis no sistema, rastreando qual \textit{User} realizou a ação, sobre qual entidade, IP e timestamp.
\end{itemize}

\section*{Planejamento Técnico do Sistema}

\begin{tabular}{@{}p{4.5cm}p{2.5cm}p{8cm}@{}}
\toprule
\textbf{Membro} & \textbf{Módulo} & \textbf{Responsabilidade} \\
\midrule
Miguel Angelo Dufloth Filho     & Auth                & Model User customizado, JWT, permissões por perfil, seed de usuários. \\
Miguel Angel Balladares Huertas & Equipamentos        & CRUD completo de equipamentos. \\
Luis Fernando Pereira           & Logs de Auditoria   & Model LogAuditoria e função \texttt{registrar\_log()}. \\
Leonardo Lotério de Lima        & Reservas e Operação & Fluxo completo de reservas, aprovação, recusa e devolução. \\
Lucas Honorato dos Santos       & Documentação        & README, \texttt{.env.example}, relatório técnico, plano de incidente. \\
\bottomrule
\end{tabular}

\section*{Evidência de Evolução do Repositório}

O repositório público encontra-se em \url{https://github.com/migueldufloth/lendify}.
O primeiro commit foi realizado em 02/06/2026, com a estrutura inicial do projeto,
README e \texttt{.gitignore}. O módulo de autenticação foi desenvolvido na branch
\texttt{feature/auth}, com PR número 1 aberto para revisão em direção à branch
\texttt{dev}. A migration inicial foi gerada e commitada, e o comando
\texttt{seed\_users} está disponível para criação dos usuários de teste.

\section*{Ativos e Dados Sensíveis}

\begin{tabular}{@{}p{3cm}p{2.5cm}p{3cm}p{5.5cm}@{}}
\toprule
\textbf{Ativo} & \textbf{Classificação} & \textbf{Onde fica} & \textbf{Risco principal} \\
\midrule
Senhas dos usuários & Restrito & Banco de dados (hash) & Acesso indevido e personificação em caso de vazamento. \\
Tokens JWT & Confidencial & Memória do Browser & Sequestro de sessão (\textit{Session Hijacking}). \\
Histórico de Reservas & Interno & Banco de dados & Vazamento de quem está portando equipamentos valiosos. \\
Logs de Auditoria & Interno & Banco de dados & Exclusão ou alteração de logs para apagar rastros de ações maliciosas. \\
\bottomrule
\end{tabular}

\end{document}
